\documentclass[11pt]{article}

\usepackage[a4paper,margin=1in]{geometry}
\usepackage{amsmath,amssymb}
\usepackage{booktabs}
\usepackage{enumitem}
\usepackage{algorithm}
\usepackage{algpseudocode}
\usepackage[hidelinks]{hyperref}

\newcommand{\R}{\mathbb{R}}
\newcommand{\refcell}{\widehat{K}}
\newcommand{\xhat}{\boldsymbol{\xi}}
\newcommand{\flop}{\mathrm{flop}}
\newcommand{\gradhat}{\widehat{\nabla}}

\title{Sampling the residual-based merit in one pass:\\
       a generated line-search kernel whose cost is independent\\
       of the number of trial step lengths}
\author{Patrick Zulian\thanks{Affiliation to be completed.}}
\date{\today}

\begin{document}
\maketitle

\begin{abstract}
A globalized Newton solver evaluates a merit function at several trial step
lengths per iteration.  When the merit is the squared residual norm
$\tfrac12\|R(x+\alpha h)\|^2$ --- the only merit available to a formulation that
is not the gradient of a potential --- the obvious implementation assembles the
whole residual once per trial step, so a search over $m$ step lengths costs $m$
assemblies.  Measured on a Mooney--Rivlin Kelvin--Voigt solid at $512{,}000$
elements on one Grace node at $72$ threads, $m=12$ costs $11.8\times$ one
assembly, and only $6\%$ of a merit evaluation is independent of the step
length.

We reorganize that evaluation in two steps.  First, a material written as
several units --- here a stored energy for the elastic response and a weak
residual for the viscous one --- has no single form to contract, because a
squared nodal residual may not be taken until every unit has contributed.  Both
units are $1$-forms on the same fields, so they sum, and the material's whole
residual becomes one kernel: one traversal in place of two, worth $1.57\times$
the element rate and $36\%$ of the time on the same node.  Second, the $m$ trial
steps share everything upstream of the constitutive law, and what they share can
be hoisted out of the trial-step loop exactly rather than approximately, because
the state enters the element kernel through a quantity that is affine in
$\alpha$.

The resulting kernel is a single pass over the mesh producing all $m$ merit
values, and it materializes no per-step residual vector: the contraction closes
inside a patch of elements, where the node sums are complete, so each trial step
leaves the patch as one number and a thread's entire output is a buffer of $m$
scalars.  What this is worth is bounded by the constitutive evaluation, which
repeats per trial step by construction.  The same reorganization is already in
place for the energy merit, and measured on the same node it costs $5.7\times$
one step at $m=12$ for a Mooney--Rivlin-class material and $3.7\times$ for
linear elasticity: the step-independent fraction rises from $6\%$ to $57\%$.
That is the quantity this work moves, and we report it in preference to the
ratio, because a material whose stress evaluation dominates pays for every step
it samples and no arrangement of the loops can make the cost flat.
\end{abstract}

\section{Notation}

Let $R:\R^n\to\R^n$ be the assembled nonlinear residual of the discrete problem,
$x\in\R^n$ the current iterate and $h\in\R^n$ the search direction.  A line
search asks for the merit
\begin{equation}
  \phi(\alpha) \;=\; \tfrac12\,\bigl\|R(x+\alpha h)\bigr\|_2^2
  \;=\; \tfrac12 \sum_{i=1}^{n} R_i(x+\alpha h)^2
  \label{eq:merit}
\end{equation}
at a prescribed set of trial step lengths $\{\alpha_k\}_{k=1}^{m}$, with $m=12$
for \texttt{hyperelasticity\_bdf2}, $m=13$ for \texttt{MaMAL} and $m=20$ by
default for \texttt{bench\_hyperelasticity}.

The residual is a sum over the operators registered on the
\texttt{Function},
\begin{equation}
  R(y) \;=\; \sum_{p\in\mathcal{P}} R_p(y),
  \label{eq:opsum}
\end{equation}
and each operator assembles element-wise from a local kernel $r^e$,
\begin{equation}
  R_p(y) \;=\; \sum_{e\in\mathcal{E}_p} P_e^{\!\top}\, r_p^e\bigl(P_e\,y\bigr),
  \label{eq:assembly}
\end{equation}
where $P_e\in\R^{N_s N_c\times n}$ is the gather of element $e$'s $N_s$ nodes and
$N_c$ field components.

\section{What is not additive, and what is}

The energy merit and the residual merit differ in exactly one respect, and it
governs the whole design.  An energy or potential is a sum of scalars, so it is
additive over operators and over elements:
\begin{equation}
  E(y) = \sum_{p} E_p(y) = \sum_{p}\sum_{e} E_p^e\bigl(P_e y\bigr).
\end{equation}
Each element contributes a number and the numbers are summed.  The residual merit
is not of that form.  Squaring happens \emph{after} the sum over operators and
after the sum over elements, and squaring does not commute with either:
\begin{equation}
  \tfrac12\Bigl\|\sum_{p} R_p\Bigr\|^2 \;\neq\; \sum_{p}\tfrac12\bigl\|R_p\bigr\|^2,
  \qquad
  \tfrac12\Bigl\|\sum_{e} P_e^{\!\top} r^e\Bigr\|^2 \;\neq\; \sum_{e}\tfrac12\bigl\|r^e\bigr\|^2 .
  \label{eq:nonadditive}
\end{equation}
Equation~\eqref{eq:nonadditive} is why \eqref{eq:merit} cannot be evaluated by an
element-local reduction the way an energy can, and why no operator can compute its
own share of it.  A node's residual is a sum over every element incident on that
node, and the square may be taken only once that sum is complete.

The consequence for the interface is that \emph{one} operator performs the
contraction, and it must be handed everything the other operators contribute.
We therefore split $\mathcal{P}$ into the operator $M$ whose residual moves with
the state and the remainder $\mathcal{A}=\mathcal{P}\setminus\{M\}$ whose residuals
do not:
\begin{equation}
  R(x+\alpha h) \;=\; \underbrace{\sum_{a\in\mathcal{A}} R_a}_{\textstyle g}
  \;+\; R_M(x+\alpha h),
  \qquad
  \frac{\partial R_a}{\partial y} \equiv 0 \quad\forall a\in\mathcal{A}.
  \label{eq:split}
\end{equation}
The vector $g\in\R^n$ --- the \emph{accumulator} --- collects the Neumann
tractions, the body force and the constraint contributions, signs included,
exactly as \texttt{Function::gradient} builds them.  It is assembled once per
Newton step rather than $m$ times, which is already correct today only by
accident: those operators are re-assembled per trial step and return the same
numbers.  With \eqref{eq:split} the merit becomes
\begin{equation}
  \boxed{\;
  \phi(\alpha_k) \;=\; \tfrac12 \sum_{i=1}^{n}
  \Bigl( g_i + \bigl[R_M(x+\alpha_k h)\bigr]_i \Bigr)^{2}
  \;}
  \label{eq:contract}
\end{equation}
and the single operator $M$ owns its evaluation for all $k$ at once.

\section{One material, several units, one form}
\label{sec:combined}

Equation~\eqref{eq:contract} assumes a single operator $M$ whose kernel produces
$R_M$.  A material need not be written that way.  Mooney--Rivlin Kelvin--Voigt is
two \emph{units} --- a stored-energy density for the elastic response and a weak
residual for the viscous one --- and each publishes its own kernels, which the
operator runs one after the other into the same vector.  Neither can contract the
merit, for the reason of \eqref{eq:nonadditive} applied one level up: after the
first traversal no node is finished, so nothing may be squared.

The resolution is not to fuse two kernels in one traversal, which would mean
carrying two element-kernel calling conventions through a single loop nest.  It
is that both units are $1$-forms on the same fields, and $1$-forms add.  Writing
$c^{(u)}_{i}$ and $\mathbf{c}^{(u)}_{i}$ for the coefficients unit $u$ contracts
against the test value and the test gradient of component $i$,
\begin{equation}
  R_i(v) \;=\; \sum_{u} \Bigl( c^{(u)}_i\,v_i
  \;+\; \mathbf{c}^{(u)}_i\!\cdot\!\nabla v_i \Bigr),
  \label{eq:formsum}
\end{equation}
so the material's whole residual is one form whose coefficients are the
componentwise sums.  A residual unit publishes its coefficients directly.  An
energy unit publishes the flux against its staged variable, and the change of
basis is the framework's own statement rather than an assumption: the staged
variable carries the expression defining it, $F = I + \nabla u$ for a deformation
gradient, so
\begin{equation}
  \mathbf{c}_i \;=\; \frac{\partial W}{\partial F_{i\cdot}}
  \Bigl|_{F = I + \nabla u}
  \;=\; \frac{\partial W}{\partial (\nabla u)_{i\cdot}} ,
\end{equation}
which was checked entry by entry against differentiating $W$ directly and agrees
exactly in rational arithmetic.

This is what makes the sampled merit possible at all, and it is worth something
on its own.  The combined form is emitted as one ordinary residual unit, so one
traversal computes what the two computed between them.  Measured on HEX8 at
$72$ threads on a Grace node, against the two units run the way the operator runs
them:

\begin{center}
\begin{tabular}{rrrrr}
\toprule
elements & dofs & two units & combined & ratio \\
         &      & [MElem/s] & [MElem/s] & \\
\midrule
 32{,}768 &   107{,}811 & 54.2 & 88.8 & 0.610 \\
110{,}592 &   352{,}947 & 60.6 & 95.8 & 0.632 \\
262{,}144 &   823{,}875 & 61.4 & 96.1 & 0.639 \\
512{,}000 & 1{,}594{,}323 & 62.5 & 97.9 & 0.639 \\
\bottomrule
\end{tabular}
\end{center}

Throughput flattens by $262{,}144$ elements, so the last two rows are saturated:
$1.57\times$ the element rate, $36\%$ of the time removed, with the two residuals
agreeing to $8\times10^{-16}$ relative.  The saving is the second traversal's
gather, geometry and trial-function accumulation --- the same quantities
Section~\ref{sec:affine} hoists out of the trial-step loop, here hoisted out of
the unit loop.

The combination is refused where it would be a second path to an existing
answer: a material already written as one unit \emph{is} its own total residual.
It is also refused for mixed-order systems, whose fields do not share a shape
count, because the lowered components carry no record of which element each came
from --- a Taylor--Hood material keeps its kernels and falls back to assembling
per trial step.

\section{The affine structure that makes one pass exact}
\label{sec:affine}

Fix an element $e$ and a quadrature point $\xhat_q$.  The element kernel reads the
state only through the value and the reference gradient of the interpolated field,
\begin{equation}
  u^e_q(y) \;=\; \sum_{j=1}^{N_s} y^e_j\,\varphi_j(\xhat_q),
  \qquad
  \gradhat u^e_q(y) \;=\; \sum_{j=1}^{N_s} y^e_j\,\gradhat\varphi_j(\xhat_q),
  \qquad y^e = P_e y .
  \label{eq:interp}
\end{equation}
Both maps $y\mapsto u^e_q$ and $y\mapsto\gradhat u^e_q$ are \emph{linear}.  With
$y = x + \alpha h$ this gives, exactly and for every $\alpha$,
\begin{equation}
  u^e_q(x+\alpha h) = u^e_q(x) + \alpha\, u^e_q(h),
  \qquad
  \gradhat u^e_q(x+\alpha h) = \gradhat u^e_q(x) + \alpha\, \gradhat u^e_q(h).
  \label{eq:affine}
\end{equation}
The right-hand sides of \eqref{eq:affine} are two accumulations over the $N_s$
shape functions, independent of $\alpha$.  The left-hand side is what the
constitutive law consumes.  So the sum over shape functions in \eqref{eq:interp}
--- the expensive part, $O(N_s\,d)$ per quadrature point per field --- is performed
\emph{twice} in total, not $2m$ times, and each trial step recovers its own state
with $d$ fused multiply--adds.

The same argument carries through the geometric map, which is also linear:
\begin{equation}
  \nabla u^e_q \;=\; J_e^{-\top}\,\gradhat u^e_q
  \quad\Longrightarrow\quad
  \nabla u^e_q(x+\alpha h) = \nabla u^e_q(x) + \alpha\,\nabla u^e_q(h),
\end{equation}
so the combination may equivalently be formed in reference space before the map or
in physical space after it.  Nothing downstream of $\nabla u^e_q$ is affine ---
the first Piola--Kirchhoff stress $P(\nabla u)$ of a Mooney--Rivlin solid is not
--- and nothing downstream is hoisted.  This is a redistribution of exact
arithmetic, not a truncation: $\phi(\alpha_k)$ is the same number the per-step
assembly computes, up to the reassociation of a sum.

\section{The kernel: a patch contracts, and emits $m$ numbers}
\label{sec:kernel}

The work unit is a \emph{patch} $p$ --- a pack of elements together with the nodes
they touch.  Its nodes divide into three classes, and the library's packed layout
already names them: the owned nodes not shared with any other pack, which we write
$\mathcal{N}_p^{\circ}$; the owned nodes that are shared; and the ghosts.  Only the
first class has all of its incident elements inside the patch, so only there is the
sum in \eqref{eq:assembly} complete without leaving it:
\begin{equation}
  \bigl[R_M(y)\bigr]_n \;=\; \sum_{\substack{e\in p \\ e\ni n}}
  \bigl[r^e_M(P_e y)\bigr]_n ,
  \qquad n\in\mathcal{N}_p^{\circ} .
  \label{eq:patchcomplete}
\end{equation}
This is not an assumption about the partition; it is the same distinction the
existing packed kernels make when they write $\mathcal{N}_p^{\circ}$ with a plain
accumulation and reach for an atomic, or for a deferred ghost buffer, on everything
else.

Where \eqref{eq:patchcomplete} holds the square in \eqref{eq:contract} is taken
inside the patch and the patch reports numbers rather than vectors.  The remaining
nodes --- the interface $\Gamma$ --- carry partial sums that must meet before they
may be squared, and they go through the deterministic ghost reduction the packed
traversal already owns, with one copy per trial step.  That buffer is the one
object in the scheme that grows with both $m$ and the mesh, and it is
\emph{surface}-sized, $m\,N_c\,|\Gamma|$ against the $m\,n$ it replaces: for a pack
of $P$ elements the interface fraction falls like $P^{-1/3}$, and the pack size is
tunable.  The fast interior and the small remainder are two passes over different
node sets, not a fast path with a fallback that redoes the work.

This is what must be understood about the data layout, and it is the point of the
whole design: \textbf{no per-step residual vector is ever materialized}.  The
contraction of trial step $k$ produces one scalar.  A thread holds
\begin{itemize}[nosep]
  \item patch-local node accumulators $\rho^{(k)}\in\R^{|\mathcal{N}_p| N_c}$,
    $k=1,\dots,m$, thread-private and sized to stay in cache --- which is what the
    pack size tunes; and
  \item one buffer of $m$ slots, $s\in\R^m$, accumulated over the patches the
    thread owns and reduced across threads at the end.
\end{itemize}
The only global per-degree-of-freedom array in the scheme is the accumulator $g$ of
\eqref{eq:split}, and it is \emph{read}, gathered per node for the pointwise
reduction and shared by all $m$ trial steps.  Nothing of size $m\,n$ is allocated,
written or streamed.

Algorithm~\ref{alg:kernel} is the emitted loop nest; phases marked $\star$ do not
depend on $\alpha$ and are executed once.

\begin{algorithm}
\caption{Residual merit at $m$ trial step lengths: one pass, $m$ scalars out}
\label{alg:kernel}
\begin{algorithmic}[1]
\State $s \gets 0 \in \R^{m}$ \Comment{thread-private, $m$ slots --- the only output}
\For{each patch $p$ owned by this thread}
  \State $\rho^{(k)} \gets 0$ for $k=1,\dots,m$
         \Comment{patch-local nodes, cache resident}
  \For{each element block $e\in p$ (vectorized over lanes)}
    \State $\star$ gather coordinates; form $J_e$, $\det J_e$, $\mathrm{adj}\,J_e$
    \State $\star$ gather $x^e$ and $h^e$ \Comment{$2 N_s N_c$ loads, once}
    \For{$q = 1,\dots,N_q$}
      \State $\star$ $\gradhat u_q \gets \sum_j x^e_j \gradhat\varphi_j(\xhat_q)$
             \Comment{$O(N_s d)$}
      \State $\star$ $\gradhat w_q \gets \sum_j h^e_j \gradhat\varphi_j(\xhat_q)$
             \Comment{$O(N_s d)$}
      \For{$k = 1,\dots,m$}
        \State $\gradhat u_q^{(k)} \gets \gradhat u_q + \alpha_k\,\gradhat w_q$
               \Comment{$d$ FMA}
        \State $\nabla u_q^{(k)} \gets J_e^{-\top}\gradhat u_q^{(k)}$
               \Comment{$O(d^2)$}
        \State evaluate the constitutive flux at $\nabla u_q^{(k)}$
        \State $\rho^{(k)} \mathrel{+}= w_q \det J_e \cdot (\text{flux}) \cdot
               \nabla\varphi_{\text{test}}$ \Comment{into $e$'s patch-local nodes}
      \EndFor
    \EndFor
  \EndFor
  \For{$n \in \mathcal{N}_p^{\circ}$} \Comment{complete: squared here}
    \State gather $g_n$ \Comment{$N_c$ loads, shared by all $m$ steps}
    \For{$k = 1,\dots,m$}
      \State $s_k \mathrel{+}= \tfrac12 \bigl\| g_n + \rho^{(k)}_n \bigr\|^2$
    \EndFor
  \EndFor
  \For{$n \in \Gamma \cap p$} \Comment{partial: deferred}
    \State append $\rho^{(k)}_n$ to the ghost buffer of step $k$
  \EndFor
\EndFor
\For{each interface node $n\in\Gamma$} \Comment{second pass, surface-sized}
  \For{$k = 1,\dots,m$}
    \State $s_k \mathrel{+}= \tfrac12 \bigl\| g_n + \textstyle\sum_{p\ni n}
           \rho^{(k)}_{n,p} \bigr\|^2$
  \EndFor
\EndFor
\State $\phi(\alpha_k) \gets \sum_{\text{threads}} s_k$
\end{algorithmic}
\end{algorithm}

The trial-step loop is placed \emph{outside} the lane loop and not inside it.  On a
CPU target the work item is a SIMD lane, and a lane loop containing another loop is
no longer innermost and does not vectorize; the same decision, for the same
measured reason, is already recorded for the stepped objective.  The
$\alpha$-independent quantities therefore go to lane-indexed buffers, the first
lane loop closes, and the trial-step loop opens a second one.

Two consequences of \eqref{eq:patchcomplete} are worth stating together, because
they are the same fact seen from two sides.  Only one operator may contract the
merit: a single contractor owns the node sums and can close them inside the patch,
whereas two contractors would each hold a partial node value and the design would
be forced back to assembling vectors.  And the gather of $g$ is per node, not per
node per step, so the state-independent part of the residual is read once for all
$m$ trial steps --- against $m$ full assemblies of it in the implementation this
replaces.

\section{Cost model}

Write $C_{\text{gath}}$ for the element gather and geometry, $C_{\text{acc}}$ for
one trial-function accumulation \eqref{eq:interp}, $C_{\text{geo}}$ for the
geometric map, $C_{\text{mat}}$ for the constitutive evaluation and
$C_{\text{con}}$ for the test-function contraction, all per element.  The
per-trial-step implementation costs
\begin{equation}
  T_{\text{naive}}(m) \;=\; m\,\bigl(C_{\text{gath}} + C_{\text{acc}}
  + C_{\text{geo}} + C_{\text{mat}} + C_{\text{con}}\bigr)
  \;+\; m\,T_{\mathcal{A}},
\end{equation}
where $T_{\mathcal{A}}$ is the assembly of the state-independent operators.  The
sampling kernel costs
\begin{equation}
  T_{\text{samp}}(m) \;=\; C_{\text{gath}} + 2\,C_{\text{acc}}
  \;+\; m\,\bigl(C_{\text{geo}} + C_{\text{mat}} + C_{\text{con}}\bigr)
  \;+\; T_{\mathcal{A}},
\end{equation}
with $T_{\mathcal{A}}$ appearing once by \eqref{eq:split}.  The asymptotic ratio is
\begin{equation}
  \lim_{m\to\infty}\frac{T_{\text{samp}}(m)}{T_{\text{naive}}(m)}
  \;=\;
  \frac{C_{\text{geo}} + C_{\text{mat}} + C_{\text{con}}}
       {C_{\text{gath}} + C_{\text{acc}} + C_{\text{geo}} + C_{\text{mat}} + C_{\text{con}} + T_{\mathcal{A}}/|\mathcal{E}|},
\end{equation}
The memory side is the sharper of the two, and it is where the patch-local
contraction earns its keep.  Writing $n$ for the number of degrees of freedom, the
per-trial-step implementation forms a combined iterate, assembles a residual into an
$n$-vector and reduces it, once per step: $\Theta(m\,n)$ words written and
$\Theta(m\,n)$ read, none of it reused.  The sampling kernel reads $x$, $h$ and $g$
--- $3n$ words, independent of $m$ --- and writes $m$ scalars per thread.  The $m$
patch-local accumulators $\rho^{(k)}$ are the only structure that grows with $m$,
they are thread-private, and their footprint $m\,|\mathcal{N}_p|\,N_c$ is bounded by
the pack size, which is tunable for exactly this reason.  The arithmetic above is
therefore the honest limit: nothing in the sampling kernel becomes bandwidth-bound
as $m$ grows.

and at $m=1$ the two agree up to the one extra accumulation $C_{\text{acc}}$.  The
measurement that settles the design is the line-search sweep over
$m\in\{1,2,4,8,12,20\}$ reported as $T(m)/T(1)$.  Measured at $512{,}000$
elements on one Grace node at $72$ threads with \texttt{OMP\_PROC\_BIND=true},
and fitted to $T(m)=a+bm$ to separate the step-independent work $a$ from the
per-step work $b$:

\begin{center}
\begin{tabular}{lrrrr}
\toprule
arm & $a$ [s] & $b$ [s] & $a/(a{+}b)$ & $T(12)/T(1)$ \\
\midrule
linear elasticity, energy   & $1.67\times10^{-3}$ & $5.36\times10^{-4}$ & 75.7\% & 3.70 \\
Mooney-Rivlin, energy       & $1.72\times10^{-3}$ & $1.29\times10^{-3}$ & 57.1\% & 5.75 \\
Mooney-Rivlin KV, residual  & $5.90\times10^{-4}$ & $9.09\times10^{-3}$ & 6.1\%  & 11.78 \\
\bottomrule
\end{tabular}
\end{center}

The last row is the target and the middle row is what a sampled kernel of
comparable constitutive cost achieves: today $6\%$ of a residual merit
evaluation is step-independent, against $57\%$ once the sharing is in place.
The number to move is that fraction, not the ratio directly.

\section{What is built, and what is not}

The form-layer combination of Section~\ref{sec:combined} is in place: the
material's whole residual is emitted as one kernel, and it agrees with the two
units the operator runs to $8\times10^{-16}$ relative on $512{,}000$ elements.
That kernel is the one the sampled merit contracts, and the measurement above is
its floor.

The sampling kernel itself is not built.  What remains is the trial-step loop
inside the element traversal, the patch-local accumulators of
Section~\ref{sec:kernel}, and the contraction to $m$ scalars.  The traversal
those need is the packed one, and it does not yet exist for this kernel: the
residual emitter's packed path is restricted to single-field systems and is
inactive for every material in the tree, so the combined kernel currently has no
packed variant.  Extending it to vector fields is the next step and is a
prerequisite rather than an optimization, because the patch is what makes the
node sums complete and therefore what makes the contraction legal at all.

One case is refused rather than approximated.  A patch is built within one mesh
block, so a node on a block boundary has incident elements the patch never sees
and would be squared while its sum is still partial --- a merit that comes out
too small with nothing to show for it.  Inter-block patches are the general
answer; until they exist the sampled path refuses a multi-block space, and the
per-step implementation, which assembles through the ordinary operator and spans
blocks like any other assembly, continues to serve it.

\section{The interface}

Two predicates and one entry point express the split \eqref{eq:split} at the
\texttt{Op} boundary:
\begin{itemize}[nosep]
  \item \texttt{energy\_or\_potential\_based()} --- whether this operator's
    $0$-form is a potential, hence summable with others.  False for a residual
    that is the gradient of nothing.
  \item \texttt{contracts\_residual\_merit()} --- whether this operator is $M$.
    A false answer is the promise $\partial R_a/\partial y\equiv 0$ that puts the
    operator in $\mathcal{A}$.
  \item \texttt{residual\_merit\_steps(x, h, nsteps, steps, accumulator, out)} ---
    evaluates \eqref{eq:contract} for every $k$ in one pass, with
    \texttt{accumulator} carrying $g$.
\end{itemize}
Because \eqref{eq:nonadditive} admits only one contractor, \texttt{add\_operator}
keeps the contracting operator last and refuses a second one.

\end{document}
